\begin{lemma}[Finite convex combinations and compactness of the convex hull]
Let $F$ be a finite-dimensional real inner-product space and let $E\subseteq F$.
Write $\operatorname{conv}E$ for the convex hull of $E$, that is, the smallest
convex subset of $F$ containing $E$. Then:
\begin{enumerate}
\item every $x\in\operatorname{conv}E$ is the mean of a finite horizontal
probability configuration supported in $E$: there are $n\ge1$, weights
$\theta_i>0$ with $\sum_{i=1}^n\theta_i=1$ and points $y_1,\dots,y_n\in E$ with
$\sum_{i=1}^n\theta_iy_i=x$, so that $(\theta,y)$ is a configuration with mean
$x$ in the sense of \noderef{IsBaryConfig} and $y_i\in E$ for every $i$;
\item if $E$ is compact then $\operatorname{conv}E$ is compact.
\end{enumerate}
\end{lemma}
\begin{proof}
Let
\[
  W=\Bigl\{\sum_{i=1}^n\theta_iy_i:n\ge1,\ y_i\in E,\ \theta_i>0,\
     \sum_{i=1}^n\theta_i=1\Bigr\}
\]
be the set of means of finite horizontal probability configurations supported in
$E$, so that (1) asserts exactly $\operatorname{conv}E\subseteq W$.

\emph{$E\subseteq W$.} For $y\in E$ take $n=1$, $\theta_1=1>0$ and $y_1=y$; then
$\sum_i\theta_i=1$ and $\sum_i\theta_iy_i=y$.

\emph{$W$ is convex.} Let $u=\sum_{i=1}^n\theta_iy_i$ and
$u'=\sum_{j=1}^{n'}\theta'_jy'_j$ be two such combinations and let
$\tau\in[0,1]$. If $\tau=0$ then $\tau u+(1-\tau)u'=u'\in W$, and if $\tau=1$
then it is $u\in W$; so assume $0<\tau<1$. Then
\[
  \tau u+(1-\tau)u'=\sum_{i=1}^n(\tau\theta_i)y_i
     +\sum_{j=1}^{n'}\bigl((1-\tau)\theta'_j\bigr)y'_j
\]
is a combination of the $n+n'\ge1$ points $y_1,\dots,y_n,y'_1,\dots,y'_{n'}$ of
$E$ with weights $\tau\theta_i>0$ and $(1-\tau)\theta'_j>0$, which sum to
$\tau\sum_i\theta_i+(1-\tau)\sum_j\theta'_j=\tau+(1-\tau)=1$. Hence
$\tau u+(1-\tau)u'\in W$, and $W$ is convex.

\emph{Proof of (1).} $W$ is a convex set containing $E$, and
$\operatorname{conv}E$ is the smallest such set; therefore
$\operatorname{conv}E\subseteq W$, which is (1).

\emph{Every member of $\operatorname{conv}E$ is a combination of at most $m+1$
points, where $m=\dim F$.} Let $u\in\operatorname{conv}E$. By (1) it has a
representation $u=\sum_{i=1}^n\theta_iy_i$ as above; choose one with $n$ least,
and suppose $n>m+1$. The $n-1$ vectors $y_i-y_n$, $1\le i\le n-1$, lie in the
$m$-dimensional space $F$ and $n-1>m$, so they are linearly dependent: there are
reals $\lambda_1,\dots,\lambda_{n-1}$, not all zero, with
$\sum_{i<n}\lambda_i(y_i-y_n)=0$. Put $\lambda_n=-\sum_{i<n}\lambda_i$, so that
\[
  \sum_{i=1}^n\lambda_iy_i=0,\qquad\sum_{i=1}^n\lambda_i=0,
\]
with not all $\lambda_i$ zero. Were every $\lambda_i\le0$, the second identity
would force them all to vanish; so $\lambda_i>0$ for some $i$. Put
$\tau=\min\{\theta_i/\lambda_i:\lambda_i>0\}$, a minimum over a nonempty finite
set, which is positive because every $\theta_i>0$; let $i_0$ attain it and put
$\theta'_i=\theta_i-\tau\lambda_i$. Then $\theta'_i\ge0$ for every $i$: for
$\lambda_i>0$ because $\tau\le\theta_i/\lambda_i$, and for $\lambda_i\le0$
because $\theta_i>0$ and $-\tau\lambda_i\ge0$. Moreover
$\sum_i\theta'_i=\sum_i\theta_i-\tau\sum_i\lambda_i=1$ and
$\sum_i\theta'_iy_i=\sum_i\theta_iy_i-\tau\sum_i\lambda_iy_i=u$, while
$\theta'_{i_0}=0$. Discarding the indices with $\theta'_i=0$ leaves a
representation of $u$ with strictly positive weights summing to $1$ and with at
least one term, since the weights sum to $1$; it has fewer than $n$ terms,
contradicting minimality. Hence $n\le m+1$.

Consequently, fixing any $y_\ast\in E$ — which exists whenever
$\operatorname{conv}E\ne\emptyset$, since by (1) a member of
$\operatorname{conv}E$ is a combination of points of $E$ and so $E\ne\emptyset$ —
and padding a representation with $n\le m+1$ terms by the terms $\theta_i=0$,
$y_i=y_\ast$ for $n<i\le m+1$,
\begin{equation}\label{eq:bcch-hull}
  \operatorname{conv}E=\Bigl\{\sum_{i=1}^{m+1}\theta_iy_i:
     \theta\in\Delta,\ y_1,\dots,y_{m+1}\in E\Bigr\},
  \qquad
  \Delta=\Bigl\{\theta\in\R^{m+1}:\theta_i\ge0,\ \sum_i\theta_i=1\Bigr\}.
\end{equation}
The padding gives the inclusion from left to right. For the other inclusion, a
sum $\sum_{i=1}^{m+1}\theta_iy_i$ with $\theta\in\Delta$ and $y_i\in E$ becomes,
after discarding the indices with $\theta_i=0$, a combination with strictly
positive weights summing to $1$ and with at least one term, hence lies in
$W$; and $W\subseteq\operatorname{conv}E$ because $\operatorname{conv}E$ is convex
and contains $E$, so it contains every finite convex combination of points of
$E$, by induction on the number of terms: one term is a point of $E$, and a
combination with $n+1$ terms and weights summing to $1$ equals
$(1-\theta_{n+1})u+\theta_{n+1}y_{n+1}$ with $u$ a combination of $n$ terms when
$\theta_{n+1}<1$, and equals $y_{n+1}\in E$ when $\theta_{n+1}=1$.

\emph{Proof of (2).} Assume $E$ compact. If $E=\emptyset$ then
$\operatorname{conv}E=\emptyset$, which is compact; so assume $E\ne\emptyset$.
$\operatorname{conv}E$ is bounded: $E$ is bounded, say $|y|\le M_E$ for $y\in E$,
and by \eqref{eq:bcch-hull} any member satisfies
$|\sum_i\theta_iy_i|\le\sum_i\theta_i|y_i|\le M_E\sum_i\theta_i=M_E$. It is
closed: let $u_k\in\operatorname{conv}E$ with $u_k\to u$ in $F$, and write
$u_k=\sum_{i=1}^{m+1}\theta^k_iy^k_i$ as in \eqref{eq:bcch-hull}. The sequence
$(\theta^k)$ lies in the bounded set $\Delta$ and each $(y^k_i)$ in the bounded
set $E$, so applying Bolzano--Weierstrass (preamble) finitely many times in
succession there is one subsequence along which $\theta^k\to\theta$ and
$y^k_i\to y_i$ for every $i$; then $\theta\in\Delta$ and $y_i\in E$ because
$\Delta$ and $E$ are closed, and
$u=\lim_k\sum_i\theta^k_iy^k_i=\sum_i\theta_iy_i$ by continuity of the finitely
many products and sums, so $u\in\operatorname{conv}E$ by
\eqref{eq:bcch-hull}. Being closed and bounded in the finite-dimensional space
$F$, the set $\operatorname{conv}E$ is compact (preamble). This is (2).
\end{proof}
